\section*{Introduction: A Reader's Map}

Parts 1 to 2 define what a business model is and how its costs behave. Parts 3 to 4 build unit economics, customer acquisition economics, retention, contribution, and operating leverage. Part 5 introduces the actual language of financial accounting: balance sheets, income statements, cash flows, revenue recognition, working capital, and the distinction between cash and profit. Part 6 explains how different industries produce different economic architectures. Part 7 provides the continuing Bayes Institute case. Part 8 is a reusable financial and business-model operating system.

This chapter is intentionally more accounting-intensive than Chapters 15 and 16. It also builds directly on Chapter 14's delivery-capacity theory and the earlier supplied essay \emph{The Theory of Pricing and Value Capture}, without confusing price acceptance with enterprise solvency. Every numerical case is hypothetical, designed to teach an internally consistent method rather than estimate Bayes's current economics. Actual reporting obligations depend on jurisdiction, entity form, transactions, and the applicable standards; obtain competent professional advice when preparing statutory accounts or tax filings.

\section*{Part I: The Commercial Machine}

\subsection*{1. The Revenue Illusion}

Consider a founder who announces, ``We made ₹10 lakh in sales this month.'' Perhaps customers paid ₹10 lakh. Perhaps the contracts were signed but nothing was collected. Perhaps the funds were prepaid for services that will consume 6 months of expensive labor. Perhaps the customers received steep rebates. Perhaps every sale costs ₹80,000 to acquire and ₹40,000 to fulfill. Without further information, the number tells us remarkably little about the quality of the business.

Entrepreneurs are surrounded by attractive numbers that measure different phenomena: leads, bookings, gross merchandise value (GMV), billings, gross revenue, cash collections, recognized revenue, gross margin, EBITDA, contribution, net income, and bank balance. Most are useful; none can substitute for the others. A founder who calls all of them ``money made'' will misallocate resources.

A strong business earns the right to repeat its sale. It produces a credible customer benefit, receives sufficient compensation, meets delivery obligations, collects cash on time, and retains resources for the next customer and future improvement. The task of this chapter is to make that system legible.

\subsection*{2. A Business Model Is a Causal Architecture, Not a Label}

``Subscription,'' ``marketplace,'' ``cohort,'' ``enterprise SaaS,'' and ``consulting'' describe elements of how firms operate; they do not fully specify why a venture can make money. A complete model answers at least 8 questions:
\begin{enumerate}
  \item \textbf{Whom} do we serve?
  \item \textbf{What job} matters?
  \item \textbf{How} is value delivered?
  \item \textbf{How} are customers reached?
  \item \textbf{Who} pays?
  \item \textbf{What} is charged for?
  \item \textbf{Which activities, assets, and partners} create the outcome?
  \item \textbf{How do the revenues, costs, cash cycles, and risks evolve} as volume increases?
\end{enumerate}

2 companies selling a ``monthly subscription'' can have entirely different profitability. 1 distributes static digital content with low incremental costs but high churn. Another provides ongoing human review with a service cost that rises with every successful customer. Both bill monthly, but they are economically different businesses.

The right way to analyze a model is through transactions and resource flows. Draw the buyer, user, payer, channel, operations, suppliers, employees, infrastructure, cash receipts, and obligations. Then trace what changes when sales double. That is more revealing than merely choosing a familiar business-model label.

\subsection*{3. Value Creation, Monetization, Capture, and Retention}

A venture can create significant customer value without being able to capture it. A free open-source tool may save companies millions; the maintainers may earn nothing. A startup can capture a high contract price but spend even more on implementation. A business can record accounting profit while the cash remains unpaid by customers. These are distinct failure modes.

\begin{itemize}
  \item \textbf{Value creation:} the improvement over feasible alternatives.
  \item \textbf{Monetization:} the mechanism through which someone pays.
  \item \textbf{Value capture:} the share of economic benefit allocated to the firm.
  \item \textbf{Retention of economic value:} what remains after the complete cost of producing, selling, financing, and renewing the service.
  \item \textbf{Cash conversion:} the timing and certainty with which the recorded earnings become usable funds.
\end{itemize}

The mechanisms are connected, but no 1 is sufficient. This is why revenue growth can mask worsening economics and why a profitable specialist service can outperform a much larger high-burn startup in founder wealth and institutional resilience.
